\begingroup
\csname tagsleft@true\endcsname
% Printed 192; principal PDF 199; Google PDF 203.
\par\clearpage
\noindent\makebox[0pt][l]{\fontsize{13}{16}\selectfont 192}\hfill{\small\textbf{NACHLASS.}}\hfill\null\par\medskip

Daraus wird zuletzt abgeleitet (in unseren Zeichen)
\[
\begin{aligned}
\log\operatorname{tang}(45^\circ+\tfrac12\omega)&=\log\operatorname{tang}(45^\circ+\tfrac12\psi)\\
&\phantom{={}}-\left(4ff+4f^4-\frac43f^6-132f^8\ldots\right)\sin\psi\\
&\phantom{={}}-\left(\frac83f^4+\frac{52}3f^6+104f^8\ldots\right)\sin3\psi\\
&\phantom{={}}-\left(\frac{88}{15}f^6+\frac{232}3f^8\ldots\right)\sin5\psi\\
&\phantom{={}}-\left(\frac{400}{21}f^8\ldots\right)\sin7\psi.
\end{aligned}
\]
[Nun ist für den Abplattungswert \(\frac{1}{302{,}7827\ldots}\) nach S.\ 183:]
{\small\setlength{\arraycolsep}{1.5pt}
\[
\begin{array}{rr@{\quad}rr@{\quad}rr@{\quad}rr}
4ff={}&10495{,}452\,5992\,\mathrm m\\
4f^4={}&4{,}325\,7581\phantom{\,\mathrm m}
 & \dfrac83 f^4={}&2{,}883\,8387\,\mathrm m\\[6pt]
-\dfrac43 f^6={}&-0{,}000\,5943\phantom{\,\mathrm m}
 & \dfrac{52}3 f^6={}&0{,}007\,7258\phantom{\,\mathrm m}
 & \dfrac{88}{15}f^6={}&0{,}002\,6149\,\mathrm m\\[6pt]
-132f^8={}&-0{,}000\,0242\phantom{\,\mathrm m}
 & 104f^8={}&0{,}000\,0191\phantom{\,\mathrm m}
 & \dfrac{232}3 f^8={}&0{,}000\,0142\phantom{\,\mathrm m}
 & \dfrac{400}{21}f^8={}&0{,}000\,0035\,\mathrm m\\
 &\overline{10499{,}777\,7388}\,\mathrm m
 & &\overline{2{,}891\,5836}\,\mathrm m
 & &\overline{0{,}002\,6291}\,\mathrm m
 & &\overline{0{,}000\,0035}\,\mathrm m\\[.6ex]
\multicolumn{8}{@{}l}{\text{[oder] in Bogensekunden [(}1\,\mathrm m=0^{\prime\prime}{,}0324\text{)]}}\\[.6ex]
 &340^{\prime\prime}{,}192\,7987\phantom{\,\mathrm m}
 & &0^{\prime\prime}{,}093\,6873\phantom{\,\mathrm m}
 & &0^{\prime\prime}{,}000\,0852\phantom{\,\mathrm m}
 & &0^{\prime\prime}{,}000\,0001.\phantom{\mathrm m}
\end{array}
\]
\par}
[Da] man hat
\[
[\operatorname{tang}(45^\circ+\tfrac12(x+iy))=
\frac{\cos x+\sin iy}{\cos iy-\sin x}=
\frac{\cos iy+\sin x}{\cos x-\sin iy},
\]
\noindent so wird]
\[
\log\operatorname{tang}(45^\circ+\tfrac12(x+iy))
=\log\surd\frac{\cos iy+\sin x}{\cos iy-\sin x}
+i\operatorname{arc}\operatorname{tang}\frac{\sin iy}{i\cos x}.
\]
Entspricht \(\Omega\) also der Breite des Ortes, dessen Koordinaten \(x,y\) sind,
und ist \(\lambda\) dessen Länge, so hat man
\[
\begin{aligned}
\mathopen{[}\log\operatorname{tang}(45^\circ+\tfrac12\Omega)+i\lambda
  &=\log\operatorname{tang}(45^\circ+\tfrac12(x+iy))\\
  &\quad-D\sin(x+iy)-E\sin3(x+iy)-F\sin5(x+iy)\ldots,
\end{aligned}
\]
\noindent wo \(D,E,F,\ldots\) die Koeffizienten der Reihe für
\(\log\operatorname{tang}(45^\circ+\tfrac12\omega)\) sind. Sind \(D,E,F,\ldots\)
in Bogensekunden gegeben, so ist mithin, mit \(\rho=\tfrac{1}{206265}\) und für
Briggs'sche Logarithmen mit dem Modul \(k\):
\[
\begin{aligned}
\lambda&=\frac1\rho\operatorname{arc}\operatorname{tang}\frac{\sin iy}{i\cos x}
-D\cos x\frac{\sin iy}{i}-E\cos3x\frac{\sin3iy}{i}-\ldots\\
\log\operatorname{tang}(45^\circ+\tfrac12\Omega)
&=\log\surd\frac{\cos iy+\sin x}{\cos iy-\sin x}
-k\rho(D\sin x\cos iy+E\sin3x\cos3iy+\ldots).
\end{aligned}
\]

% Printed 193; principal PDF 200; Google PDF 204.
\par\clearpage
\noindent\null\hfill{\small\textbf{KONFORME ABBILDUNG DES SPHÄROIDS IN DER EBENE.}}\hfill\makebox[0pt][r]{\fontsize{13}{16}\selectfont 193}\par\medskip

Mit den angegebenen Werten für \(D,E,F,\ldots\) erhält man daher:]
\[
\begin{aligned}
\lambda=\frac1\rho\operatorname{arc}\operatorname{tang}\frac{\sin iy}{i\cos x}
&-340^{\prime\prime}{,}192\,7987\cos x\frac{\sin iy}{i}\\
&-\phantom{34}0^{\prime\prime}{,}093\,6873\cos3x\frac{\sin3iy}{i}\\
&-\phantom{34}0^{\prime\prime}{,}000\,0852\cos5x\frac{\sin5iy}{i}\\
&-\phantom{34}0^{\prime\prime}{,}000\,0001\cos7x\frac{\sin7iy}{i}\\
&\quad\text{etc.}
\end{aligned}
\]
\[
\begin{aligned}
\log\operatorname{tang}(45^\circ+\tfrac12\Omega)
&=\log\surd\frac{\cos iy+\sin x}{\cos iy-\sin x}\\
&\phantom{={}}-0{,}000\,716\,2824\sin x\cos iy\\
&\phantom{={}}-0{,}000\,000\,1973\sin3x\cos3iy\\
&\phantom{={}}-0{,}000\,000\,0002\sin5x\cos5iy\\
&\qquad\qquad\text{etc.}
\end{aligned}
\]
[Mithilfe der Tabelle auf S.\ 184 ergibt sich dann zu \(\Omega\) das zugehörige \(\Phi\).]

\begin{center}[32.]\end{center}

[Die Darstellung der Oberfläche des Sphäroids in der Ebene wird] durch folgende zwei Formeln ausgedrückt:
\begin{equation*}
\operatorname{cotang}\tfrac12 p\left(\frac{1-e\cos p}{1+e\cos p}\right)^{\frac12 e}
\cdot(\cos\lambda+i\sin\lambda)
=\operatorname{cotang}\tfrac12 P\left(\frac{1-e\cos P}{1+e\cos P}\right)^{\frac12 e}
\tag*{1)}
\end{equation*}
\begin{equation*}
x+iy=\int\frac{a(1-ee)\,\mathrm dP}{(1-ee\cos P^2)^{\frac32}},
\tag*{2)}
\end{equation*}
\noindent wo \(a\) und \(a\surd(1-ee)\) die beiden Halbachsen des Ellipsoids sind
[und \(p=90^\circ-\Phi\) gesetzt ist.

Setzt man
\[
\operatorname{cotang}\tfrac12 p\left(\frac{1-e\cos p}{1+e\cos p}\right)^{\frac12 e}
=\operatorname{cotang}\tfrac12 q
\]
\noindent und]
\[
\operatorname{cotang}\tfrac12 q(\cos\lambda+i\sin\lambda)=\operatorname{cotang}\tfrac12 Q,
\]
\noindent [sowie]
\[
Q=R+iS,
\]
\noindent\textsc{ix.}\hfill 25\hspace*{.12\linewidth}\par

% Printed 194; principal PDF 201; Google PDF 205.
\par\clearpage
\noindent\makebox[0pt][l]{\fontsize{13}{16}\selectfont 194}\hfill{\small\textbf{NACHLASS.\quad KONFORME ABBILDUNG DES SPHÄROIDS IN DER EBENE.}}\hfill\null\par\medskip

\noindent [so wird]
\[
\begin{aligned}
\operatorname{tang}\tfrac12 q(\cos\lambda-i\sin\lambda)&=\operatorname{tang}\tfrac12(R+iS)\\
\operatorname{tang}\tfrac12 q(\cos\lambda+i\sin\lambda)&=\operatorname{tang}\tfrac12(R-iS).
\end{aligned}
\]
[Hieraus folgt:
\[
\begin{alignedat}{2}
2\operatorname{tang}\tfrac12 q\cos\lambda&=\frac{\sin R}{\tfrac12(\cos R+\cos iS)},\qquad&
-2i\operatorname{tang}\tfrac12 q\sin\lambda&=\frac{\sin iS}{\tfrac12(\cos R+\cos iS)},\\
1-\operatorname{tang}\tfrac12 q^2&=\frac{\cos R}{\tfrac12(\cos R+\cos iS)},&
1+\operatorname{tang}\tfrac12 q^2&=\frac{\cos iS}{\tfrac12(\cos R+\cos iS)};
\end{alignedat}
\]
\noindent mithin]
\[
\begin{aligned}
\operatorname{tang} R&=\frac{2\cos\lambda\operatorname{tang}\tfrac12 q}{1-\operatorname{tang}\tfrac12 q^2}
=\cos\lambda\operatorname{tang} q\\
\operatorname{tang} iS&=-\frac{2i\sin\lambda\operatorname{tang}\tfrac12 q}{1+\operatorname{tang}\tfrac12 q^2}
=-i\sin\lambda\sin q
\end{aligned}
\]
\noindent [oder]
\[
S=\tfrac12\log\frac{1-\sin\lambda\sin q}{1+\sin\lambda\sin q}.
\]
\begin{center}\rule{.18\linewidth}{.4pt}\end{center}
\small
\setlength{\abovedisplayskip}{5pt}
\setlength{\belowdisplayskip}{5pt}
\setlength{\abovedisplayshortskip}{3pt}
\setlength{\belowdisplayshortskip}{3pt}

% Printed 195; principal PDF 202; Google PDF 206.
\par\clearpage

\thispagestyle{empty}
\vspace*{1.6\baselineskip}
\begin{center}\large\textbf{BEMERKUNGEN.}\end{center}
{\small
Eine zusammenhängende Darstellung der konformen Abbildung des Sphäroids in der Ebene,
die als Projektionsmethode der hannoverschen Gradmessung und Landesvermessung diente,
hat sich im Nachlass nicht gefunden. Die vorstehenden Entwicklungen sind mithilfe der
Aufzeichnungen in zwei Notizbüchern und einer Reihe einzelner Blätter zusammengestellt
worden, die \textsc{Gauss} vielleicht als Rechenpapiere bei der Ableitung seiner Formeln benutzte.

Das Notizbuch »Den astronomischen Wissenschaften gewidmet« enthält die Formeln zur
Berechnung der geografischen Koordinaten, der Meridiankonvergenz und des
Vergrößerungsverhältnisses aus den ebenen rechtwinkligen Koordinaten. Direkt entnommen
sind ihm die Artikel [4], [9], [10], [11], ferner die Artikel über die Azimutreduktion [15]
mit der Tabelle für \(\log q\), [16], [17] und die über die Transformation der Koordinaten
[18], [19], [20]. Aus verschiedenen Stellen dieses Notizbuchs sind die Artikel [7], [13]
und [14] zusammengesetzt.

Ein zweites Notizbuch mit dem Titel »Kleine Aufsätze aus verschiedenen Theilen der
Mathematik« lieferte die Artikel [5], [26], [27] und [31]. Für [27] wurde außerdem noch
ein einzelnes Blatt benutzt.

Die Artikel [1], [2], [6], [8], [21] und [22] sind nach Aufzeichnungen des ersten
Notizbuchs und denen einzelner Blätter zusammengestellt worden.

Einzelnen Blättern entnommen sind [3] mit der dazugehörigen Tabelle, [12], [23], [29]
und [32] sowie die zweite Tabelle zu [15], die hier vollständig mitgeteilt ist, weil keine
andere vorhanden ist, die auf einem anderen Abplattungswert als \(1:302{,}68\) beruht.
Diese Tabelle befindet sich in einem besonderen Heftchen. Die Bezeichnungen der einzelnen
Spalten sind hier wie auch bei der ersten Tabelle zu [15] hinzugefügt worden.
\textsc{Gauss} hatte über die erste Spalte geschrieben: Corrige \(-26{,}264\), was,
wie bereits erwähnt, einer Verkleinerung der Ausgangsbreite um \(0^{\prime\prime}{,}85\)
entspricht. Die Tabelle zu [28], von der oben nur ein Auszug gegeben ist --- sie
erstreckt sich im Original von \(\psi=5000000\,\mathrm m\) bis \(\psi=6500000\,\mathrm m\) ---
und die Tabelle unter [30] befinden sich auf demselben Blatt.

Die Artikel [24], [25] und [28] schließlich sind nach Aufzeichnungen auf verschiedenen
Blättern zusammengesetzt.

\medskip
In den einzelnen Artikeln ist eine Anzahl von Änderungen vorgenommen worden.

\medskip
\noindent Bei [2]:\quad Anstelle der letzten Formel für \(c\) hat \textsc{Gauss} auf einem einzelnen Blatt
\[
c=\lambda\sin\varphi-
\frac{(1-ee\sin\varphi^2)^{\frac32}\operatorname{tang}\varphi}{6a^3(1-ee)^2}
\left\{1-ee+ee(3+ee)\cos\varphi^2+4e^4\cos\varphi^4\right\}y^3.
\]
\noindent »\quad [4]:\quad Bei der Formel für \(c\) fehlt bei ihm im Ausdruck für \(\log\operatorname{hyp} C'\) das Glied
\[
-\tfrac18\frac{1-ee\sin\varphi^2}{aa}yy.
\]
\noindent »\quad [7]:\quad Anstelle der ersten Formel steht im Original
\[
N+i\log\mathrm{Conv.}=\log f'(\xi)-\log f'(x+iy),\qquad N=\log n,
\]
\noindent\hspace*{3.5em}was aber wohl nur ein Schreibfehler ist.
\par}
\noindent\hfill \(25^{*}\)\hspace*{.12\linewidth}\par

% Printed 196; principal PDF 203; Google PDF 207.
\par\clearpage
\noindent\makebox[0pt][l]{\fontsize{13}{16}\selectfont 196}\hfill{\small\textbf{BEMERKUNGEN.}}\hfill\null\par\medskip

{\small
\noindent Bei [8]:\quad Die letzte Formel, die sich auf einem einzelnen Blatt befindet, heißt dort
\[
n=1-\frac{\gamma}{2\alpha}yy+
\frac{\alpha\alpha\varepsilon-6\alpha\beta\delta+6\beta\beta\gamma}{24\alpha^3}y^4.
\]
\noindent »\quad [9]:\quad In der Formel für \(\log n\) fehlt bei \textsc{Gauss} im Koeffizienten von \(y^6\) das Glied \(\tfrac1{720}A^{\mathrm{IV}}\).

\noindent »\quad [11]:\quad Im Original steht anstelle der angegebenen Formeln für \(x\) und \(y\):
\[
\begin{aligned}
x&=\xi-\tfrac12\alpha\beta\lambda\lambda
+\tfrac1{24}(\alpha\beta^3+4\alpha\alpha\beta\gamma+\alpha^3\delta)\lambda^4\\
y&=\alpha\lambda-\tfrac16(\alpha\alpha\gamma+\alpha\beta\beta)\lambda^3.
\end{aligned}
\]
\noindent »\quad [12]:\quad Der Faktor von \(y^6\) in der Formel für \(\xi\) lautet auf einem Zettel, der mit anderen bei \textsc{Gauss}' Tod auf seinem Schreibtisch lag,
\[
\tfrac1{45}\alpha_1^5+\tfrac5{144}\alpha_1^3\beta_1
-\tfrac1{144}\alpha_1\alpha_1\gamma_1+\tfrac1{48}\alpha_1\beta_1\beta_1
+\tfrac1{72}\beta_1\gamma_1-\tfrac1{144}\alpha_1\delta_1+\tfrac1{720}\varepsilon_1.
\]
\noindent »\quad [13]:\quad Im Ausdruck für \(\mathrm g^{\mathrm V}(\mathrm F)\)
und entsprechend in der ersten Formel für \(y\) beim Faktor von \(\lambda^5\),
S.\ 157, lautet der Wert in der Klammer im Original:
\[
(1-20cc+(24-57\delta)c^4+\ldots),
\]
\noindent während es heißen muss, wie oben angegeben:
\[
(1-20cc+(24-58\delta)c^4+\ldots).
\]
\noindent »\quad [15]:\quad Anstelle von \(\log\mathfrak{B}=7{,}317\,8248-30\) hat
\textsc{Gauss} \(\log\mathfrak{B}=7{,}317\,8148-30\). Infolgedessen sind in der
zweiten Tabelle des Art.\ [15], S.\ 161 und 162, die Werte für \(\log\beta\) um
eine Einheit der letzten Stelle zu erhöhen.

\noindent »\quad [16]:\quad Der Koeffizient von \(R^3\) in der Formel für \(-n^*\sin\varphi\),
S.\ 164, lautet im Original
\[
\tfrac14c'+\frac{ab'}{24N}-\frac{a'b''}{24N}
+\frac{a^{\prime3}}{48NN}+\frac{aaa'}{24NN}.
\]
\noindent \textsc{Gauss} schreibt in der Formel für \(\log n^*\), S.\ 165 oben,
die Glieder zweiter Dimension
\[
-Cxx-2Dxy-Eyy;
\]
\noindent dazu passt aber die folgende Entwicklung nicht.

\noindent »\quad [17]:\quad Für \(r\), S.\ 168, ist im Original angegeben
\[
r=R+\tfrac23ttR^3+(\tfrac32tu-\tfrac12t^4)R^4\ldots;
\]
\noindent demzufolge lautet weiter die Differentialgleichung für \(\psi\)
\[
\begin{aligned}
&\left(1+(2tt+t't')RR+(6tu-2t^4+2t'u')R^3\ldots\right)
 \frac{\partial\psi}{\partial R}\\
&\qquad+\left(-\tfrac23tt'R-(\tfrac12tu'+\tfrac32t'u+2t't^3)RR\ldots\right)
 \frac{\partial\psi}{\partial\theta}=0,
\end{aligned}
\]
\noindent woraus sich dann ergibt:
\[
\psi=\theta+\tfrac13tt'RR+\left((\tfrac12u+\tfrac23t^3)t'+\tfrac16tu'\right)R^3\ldots.
\]
\noindent »\quad [20]:\quad \textsc{Gauss} hat in der Formel für \(x+\Omega\) im Faktor von \(\omega\)
\[
1-\tfrac12(1+tt)yy+(\tfrac5{24}-\tfrac{13}{24}tt-\tfrac{15}8t^4)y^4\ldots.
\]
\noindent »\quad [23]:\quad Im Original heißt der Koeffizient von \(y^9\) in der Formel für \(x\), S.\ 175,
\[
\frac{18621\delta^4+48826\delta^5+48160\delta^6+21288\delta^7+3548\delta^8}{2835},
\]
\noindent und dementsprechend lautet der Koeffizient von \(y^{10}\) in der Formel für \(xx\):
\[
\frac{316305\delta^4+831340\delta^5+822217\delta^6+363738\delta^7+60623\delta^8}{14175};
\]
\par}

% Printed 197; principal PDF 204; Google PDF 208.
\par\clearpage
\noindent\null\hfill{\small\textbf{KONFORME ABBILDUNG DES SPHÄROIDS IN DER EBENE.}}\hfill\makebox[0pt][r]{\fontsize{13}{16}\selectfont 197}\par\medskip

{\small
\noindent In der darauffolgenden Formel für \(1-xx\) hat das Original als Koeffizienten von \(y^8\):
\[
\frac{7029\delta^4+16182\delta^5+21574\delta^6+5798\delta^7+892\delta^8}{315}.
\]
\noindent Bei [24]:\quad Der Faktor von \(\sin\varphi^5\cos\varphi\) in der zweiten Formel
für \(\psi\), S.\ 177 oben, heißt im Original:
\[
\tfrac{241}{480}e^6-\tfrac{3777}{30420}e^8.
\]
\noindent Als Ausdruck für \(\sin\varphi\cos\varphi\), S.\ 177, gibt \textsc{Gauss} an:
\[
\sin\omega\cos\omega\left(1+ee+e^4+e^6-\frac{12ee+25e^4+39e^6}{6}yy
+\frac{65e^4+258e^6}{15}y^4-\frac{146}{15}e^6y^6\right).
\]
\noindent Die Formel für \(\psi\), durch \(\omega\) ausgedrückt, S.\ 177 unten, lautet im Original:
\[
\begin{aligned}
\psi=\omega&+\sin2\omega\left(\frac18ee-\frac{13}{192}e^4+\frac{173}{1024}e^6
+\frac{29029}{368640}e^8\ldots\right)\\
&-\sin4\omega\left(\frac{47}{768}e^4-\frac{293}{3840}e^6+\frac{1220619}{884736}e^8\ldots\right)\\
&+\sin6\omega\left(\frac{979}{15360}e^6-\frac{4124089}{10321920}e^8\ldots\right)\\
&-\sin8\omega\left(\frac{12895751}{123863040}e^8\ldots\right).
\end{aligned}
\]
\noindent Bei \textsc{Gauss} heißt der Faktor von \(-\sin2\varphi\) in der letzten Formel,
die \(\omega\) als Funktion von \(\varphi\) liefert, S.\ 178,
\[
\tfrac12ee+\tfrac5{24}e^4+\tfrac3{32}e^6+\tfrac{281}{11520}e^8.
\]
\noindent »\quad [26]:\quad Im Original steht anstelle der ersten beiden Zeilen im Ausdruck
für \(\log\cos\varphi\):
\[
\begin{aligned}
\log\cos\varphi&=\log\cos\psi-6ff-63f^4-\frac{2192}3f^6-\frac{17559}2f^8\ldots\\
&\qquad+\left(6ff+42f^4+\frac{946}3f^6+2200f^8\ldots\right)\cos2\psi.
\end{aligned}
\]
\noindent »\quad [28]:\quad \textsc{Gauss} hat in der Formel, die \(\psi\) als Funktion
von \(\omega\) darstellt, S.\ 183,
\[
4{,}800\,0528\sin4\omega
\]
\[
\text{anstelle von }4{,}700\,0528\sin4\omega,
\]
\noindent und dementsprechend in der Rechnung bei [29], S.\ 188,
\(\log B=0{,}681\,2460\) anstelle von \(\log B=0{,}672\,1028\). Dadurch ist bei ihm
der Wert von \(y\), S.\ 188, um \(0{,}007\,\mathrm m\) kleiner als oben angegeben.

\noindent Der bei den Entwicklungen unter [28] benutzte Wert
\[
\log\frac{2\,.\,10^7}{\pi}ff=3{,}418\,9411.803.
\]
\noindent der auch der Tabelle, S.\ 184, die den Übergang von \(\psi\) zu \(\varphi\)
und \(\omega\) vermittelt, zugrunde liegt, ist wahrscheinlich durch einen Schreibfehler
bei \textsc{Gauss} entstanden. Auf einem einzelnen Blatt ist nämlich dem eben
angegebenen Wert von \(\log\frac{2\,.\,10^7}{\pi}ff\) beigefügt: \(n=302{,}78\).
Das trifft aber nicht zu, sondern zu diesem Wert der Abplattung gehört
\[
\log\frac{2\,.\,10^7}{\pi}ff=3{,}418\,9451.803.
\]
\par}

% Printed 198; principal PDF 205; Google PDF 209.
\par\clearpage
\noindent\makebox[0pt][l]{\fontsize{13}{16}\selectfont 198}\hfill{\small\textbf{BEMERKUNGEN.}}\hfill\null\par\medskip

{\small
\noindent Bei [31]:\quad In der zweiten Formel für \(\log\operatorname{tang}\tfrac12 p\), S.\ 191, steht im Original:
\[
\begin{array}{cccc@{\quad}ccc@{\quad}rcl}
\text{im}&\text{Koeffizienten}&\text{von}&-\cos q&
 \text{an}&\text{Stelle}&\text{von}&\tfrac{11}{48}a^3&:&\tfrac{11}{24}a^3\\
\text{»}&\text{»}&\text{»}&-\cos3q&
 \text{»}&\text{»}&\text{»}&\tfrac3{32}a^3&:&\tfrac3{16}a^3\\
\text{»}&\text{»}&\text{»}&-\cos5q&
 \text{»}&\text{»}&\text{»}&\tfrac1{96}a^3&:&\tfrac1{48}a^3.
\end{array}
\]
\noindent Infolgedessen heißt es bei \textsc{Gauss}:
\[
\begin{aligned}
\log\operatorname{tang}(45^\circ+\tfrac12\omega)
 &=\log\operatorname{tang}(45^\circ+\tfrac12\psi)
 -\left(4ff+4f^4+\frac{40}3f^6-22f^8\ldots\right)\sin\psi\\
 &\quad-\left(\frac83f^4+\frac{34}3f^6+59f^8\ldots\right)\sin3\psi
 -\left(\frac{98}{15}f^6+\frac{247}3f^8\ldots\right)\sin5\psi\\
 &\quad-\left(\frac{400}{21}f^8\ldots\right)\sin7\psi.
\end{aligned}
\]
\noindent Als numerische Werte der Koeffizienten dieser Reihe gibt \textsc{Gauss} an:
\[
\begin{array}{r@{\quad}l@{\quad}l@{\quad}l}
10494{,}784\,2962\,\mathrm m;&2{,}888\,891\,3195\,\mathrm m;&
0{,}002\,927\,170\,360\,\mathrm m;&0{,}000\,003\,499\,170\,05\,\mathrm m;\\
340^{\prime\prime}{,}031\,0112;&0^{\prime\prime}{,}093\,600\,078\,75;&
0^{\prime\prime}{,}000\,094\,840\,32;&0^{\prime\prime}{,}000\,000\,113\,373\,11;
\end{array}
\]
\noindent und als Schlussformeln:
\[
\begin{aligned}
\lambda={}&\operatorname{arc}\operatorname{tang}\frac{\sin iy}{i\cos x}\\
&-340^{\prime\prime}{,}031\,0112\cos x\frac{\sin iy}{i}
 -0^{\prime\prime}{,}093\,6001\cos3x\frac{\sin3iy}{i}\\
&-0^{\prime\prime}{,}000\,0948\cos5x\frac{\sin5iy}{i}
 -0^{\prime\prime}{,}000\,0001\cos7x\frac{\sin7iy}{i},
\end{aligned}
\]
\[
\begin{aligned}
\log\operatorname{tang}(45^\circ+\tfrac12\Omega)
={}&\log\surd\frac{\cos iy+\sin x}{\cos iy-\sin x}\\
&-0{,}000\,715\,9403\sin x\cos iy
 -0{,}000\,000\,1971\sin3x\cos3iy\\
&-0{,}000\,000\,0002\sin5x\cos5iy.
\end{aligned}
\]
\noindent Bei [32]:\quad Im Ausdruck für \(S\), S.\ 194, fehlt im Original der Faktor \(\tfrac12\).
\begin{center}\rule{.24\linewidth}{.4pt}\end{center}

Zu den Formeln selbst möge noch Folgendes bemerkt werden.

Die Grundformeln der \textsc{Gauss}'schen Entwicklung, Art.\ [1], erhält man, indem man
die Ebene konform auf das Sphäroid überträgt. Die Aufgabe ist also die Umkehrung der in
Art.\ 12 der »Allgemeinen Auflösung der Aufgabe: Die Theile einer gegebnen Fläche auf einer
andern gegebnen Fläche so abzubilden, dass die Abbildung dem Abgebildeten in den kleinsten
Theilen ähnlich wird«, Bd.\ IV, S.\ 205 u.\ f., behandelten Übertragung. Es tritt aber hier
die Bedingung hinzu, dass die Strecken auf der Abszissenachse in der Ebene den entsprechenden
Bögen eines bestimmten Meridians, des Hauptmeridians, gleich sein sollen.
Sind \(x,y\) die rechtwinkligen Koordinaten eines Punktes in der Ebene, \(\Phi,\lambda\) die
geografischen Koordinaten des entsprechenden Punktes auf dem Sphäroid, so ist im Anschluss
an Art.\ 12:
\[
\omega=\mathrm dx^2+\mathrm dy^2,
\]
\noindent und die Differentialgleichung \(\omega=0\) gibt \(x\pm iy=\mathrm{const.}\)

Ferner ist
\[
\begin{aligned}
\Omega&=\frac{aa(1-ee)^2}{(1-ee\sin\Phi^2)^3}\,\mathrm d\Phi^2
 +\frac{aa\cos\Phi^2}{1-ee\sin\Phi^2}\,\mathrm d\lambda^2\\
&=\frac{aa\cos\Phi^2}{1-ee\sin\Phi^2}
\left\{\left(\frac{1-ee}{(1-ee\sin\Phi^2)\cos\Phi}\right)^2
 \mathrm d\Phi^2+\mathrm d\lambda^2\right\};
\end{aligned}
\]
\noindent die Differentialgleichung \(\Omega=0\) ergibt mithin
\[
\frac{1-ee}{(1-ee\sin\Phi^2)\cos\Phi}\,\mathrm d\Phi\pm i\,\mathrm d\lambda=0,
\]
\par}

% Printed 199; principal PDF 206; Google PDF 210.
\par\clearpage
\noindent\null\hfill{\small\textbf{KONFORME ABBILDUNG DES SPHÄROIDS IN DER EBENE.}}\hfill\makebox[0pt][r]{\fontsize{13}{16}\selectfont 199}\par\medskip

{\small
\noindent oder wenn
\begin{equation*}
\frac{(1-ee\sin\Phi^2)\cos\Phi}{1-ee}=\Theta(\Phi)
\tag*{1)}
\end{equation*}
\noindent gesetzt wird,
\[
\frac{\mathrm d\Phi}{\Theta(\Phi)}\pm i\,\mathrm d\lambda=0,
\]
\noindent deren Integral
\[
\int\frac{\mathrm d\Phi}{\Theta(\Phi)}\pm i\lambda=\mathrm{const.},
\]
\noindent oder
\begin{equation*}
\mathrm F(\Phi)=\int\frac{\mathrm d\Phi}{\Theta(\Phi)}
=\log\operatorname{hyp}\operatorname{tang}(45^\circ+\tfrac12\Phi)
+\tfrac12 e\log\operatorname{hyp}\frac{1-e\sin\Phi}{1+e\sin\Phi}
\tag*{2)}
\end{equation*}
\noindent gesetzt:
\[
\mathrm F(\Phi)\pm i\lambda=\mathrm{const.}
\]
\noindent ist.

Zu der Breite \(\Phi\) soll der Meridianbogen \(\xi\) gehören, vom Äquator ab gerechnet,
sodass also auch
\begin{equation*}
\mathrm F(\Phi)=f(\xi)=\int\frac{\mathrm d\Phi}{\Theta(\Phi)}
\tag*{3)}
\end{equation*}
\noindent ist.

Bedeutet \(\mathfrak F\) eine willkürliche Funktion, so wird mithin die Ebene konform
auf das Ellipsoid übertragen durch die Gleichung
\[
f(\xi)+i\lambda=\mathfrak F(x+iy).
\]
Nun soll aber jede Strecke auf der Abszissenachse gleich dem entsprechenden Meridianbogen
sein. Für \(\lambda=0\) muss also \(\xi=x\) und \(y=0\) sein. Folglich wird
\[
f(x)=\mathfrak F(x).
\]
\noindent Mithin lautet die Übertragungsgleichung:
\begin{equation*}
f(\xi)+i\lambda=\mathrm F(\Phi)+i\lambda=f(x+iy).
\tag*{4)}
\end{equation*}
\begin{center}\rule{.24\linewidth}{.4pt}\end{center}

Die Ableitung der Formeln für das Vergrößerungsverhältnis, Art.\ [6] bis [10],
kann wie folgt geschehen.

Da
\[
\begin{gathered}
\mathrm df(\xi)=\frac{\mathrm d\Phi}{\Theta(\Phi)}
=\frac{(1-ee)\,\mathrm d\Phi}{\cos\Phi(1-ee\sin\Phi^2)}\\
f'(\xi)=\frac{1-ee}{\cos\Phi(1-ee\sin\Phi^2)}
\,.\,\frac{(1-ee\sin\Phi^2)^{\frac32}}{a(1-ee)}
=\frac{\surd(1-ee\sin\Phi^2)}{a\cos\Phi},
\end{gathered}
\]
\noindent wobei \(\xi\) mit \(\Phi\) wächst, also nach Norden positiv angenommen ist, so wird
\[
\begin{aligned}
\frac1{nn}=\frac{\Omega}{\omega}
 &=\left(\frac1{f'(\xi)}\right)^2
 \frac{(\mathrm df(\xi))^2+\mathrm d\lambda^2}{\mathrm dx^2+\mathrm dy^2}
 =\left(\frac1{f'(\xi)}\right)^2
 \frac{\mathrm d(f(\xi)+i\lambda)}{\mathrm d(x+iy)}
 \,.\,\frac{\mathrm d(f(\xi)-i\lambda)}{\mathrm d(x-iy)}\\
 &=\left(\frac1{f'(\xi)}\right)^2
 \left\{\left(\frac{\partial\lambda}{\partial x}\right)^2+
 \left(\frac{\partial\lambda}{\partial y}\right)^2\right\}.
\end{aligned}
\]
Führt man die Meridiankonvergenz \(c\) ein, so ist zunächst
\(\operatorname{tang}c=-\dfrac{\mathrm dy}{\mathrm dx}\), wobei
\(\dfrac{\mathrm dy}{\mathrm dx}\) aus der Gleichung
\(\lambda=\mathrm{const.}\) erhalten wird. Mithin ist
\par}

% Printed 200; principal PDF 207; Google PDF 211.
\par\clearpage
\noindent\makebox[0pt][l]{\fontsize{13}{16}\selectfont 200}\hfill{\small\textbf{BEMERKUNGEN.}}\hfill\null\par\medskip

{\small
\[
\frac1{\cos c^2}=\frac{\left(\dfrac{\partial\lambda}{\partial x}\right)^2
 +\left(\dfrac{\partial\lambda}{\partial y}\right)^2}
 {\left(\dfrac{\partial\lambda}{\partial y}\right)^2},\qquad
\frac1{\sin c^2}=\frac{\left(\dfrac{\partial\lambda}{\partial x}\right)^2
 +\left(\dfrac{\partial\lambda}{\partial y}\right)^2}
 {\left(\dfrac{\partial\lambda}{\partial x}\right)^2},
\]
\noindent und daher
\begin{equation*}
\GaussOldenburgLeftBrace{\begin{aligned}
\frac{\partial\lambda}{\partial y}
 &=\phantom{-}\frac{\partial f(\xi)}{\partial x}=\frac1n f'(\xi)\cos c=u\\
\frac{\partial\lambda}{\partial x}
 &=-\frac{\partial f(\xi)}{\partial y}=\frac1n f'(\xi)\sin c=v.
\end{aligned}}
\tag*{1)}
\end{equation*}
Einfacher ergeben sich diese Formeln, wenn man bedenkt, dass in der Ebene
(vorausgesetzt, dass \(x\) nach Norden positiv ist und \(y\) mit \(\lambda\) wächst)
für das Bild eines Elements des Parallelkreises:
\begin{equation*}
\frac{n\,.\,a\cos\Phi}{\surd(1-ee\sin\Phi^2)}\,\mathrm d\lambda
=\frac n{f'(\xi)}\,\mathrm d\lambda=\mathrm dx\sin c+\mathrm dy\cos c
\tag*{2)}
\end{equation*}
\noindent ist.

Aus 1) erhält man:
\begin{equation*}
\frac1{nn}=\left(\frac u{f'(\xi)}\right)^2+\left(\frac v{f'(\xi)}\right)^2.
\tag*{3)}
\end{equation*}
Mithilfe dieser Formel hat \textsc{Gauss} in [6] das Vergrößerungsverhältnis berechnet.

Man hätte es noch etwas leichter durch die aus 3) folgende Formel
\begin{equation*}
\frac1{nn}=\left(\frac{\partial\xi}{\partial x}\right)^2+
\left(\frac{\partial\xi}{\partial y}\right)^2
\tag*{4)}
\end{equation*}
\noindent erhalten. Direkt findet man diese Formel, indem man für das Bild eines Elements
des Meridians ansetzt:
\[
\frac{n\,.\,a(1-ee)}{\surd(1-ee\sin\Phi^2)^3}\,\mathrm d\Phi
=n\,\mathrm d\xi=\mathrm dx\cos c-\mathrm dy\sin c,
\]
\noindent woraus folgt
\[
\frac{\partial\xi}{\partial x}=\frac1n\cos c,\qquad
\frac{\partial\xi}{\partial y}=-\frac1n\sin c.
\]
Nach 1) ist ferner
\[
\frac1{f'(\xi)}\left\{\frac{\partial f(\xi)}{\partial x}
+i\frac{\partial\lambda}{\partial x}\right\}
=\frac1n(\cos c+i\sin c)
\]
\noindent oder
\[
\frac{f'(x+iy)}{f'(\xi)}=\frac1n e^{ic},
\]
\noindent oder wenn
\[
\log n=N
\]
\noindent gesetzt wird,
\begin{equation*}
N-ic=\log f'(\xi)-\log f'(x+iy).
\tag*{5)}
\end{equation*}
Den hieraus folgenden Wert
\[
N=\log f'(\xi)-\text{Realteil von }\log f'(x+iy)
\]
\noindent erhält man auch sofort aus der zuerst angegebenen Formel auf der vorigen Seite.
Nach dieser wird
\[
n=\frac{f'(\xi)}{\surd\{f'(x+iy)f'(x-iy)\}}.
\]
Wird die Gleichung 5) zweimal nach \(x\) und \(y\) differenziert, so findet man wegen 4),
wenn man
\par}

% Printed 201; principal PDF 208; Google PDF 212.
\par\clearpage
\noindent\null\hfill{\small\textbf{KONFORME ABBILDUNG DES SPHÄROIDS IN DER EBENE.}}\hfill\makebox[0pt][r]{\fontsize{13}{16}\selectfont 201}\par\medskip

{\small
\noindent außerdem berücksichtigt, dass
\[
\frac{\partial\partial f(\xi)}{\partial x^2}
+\frac{\partial\partial f(\xi)}{\partial y^2}=0
\]
\noindent ist:
\begin{equation*}
\frac{\partial\partial c}{\partial x^2}
+\frac{\partial\partial c}{\partial y^2}=0
\tag*{6)}
\end{equation*}
\noindent und
\begin{equation*}
\frac{\partial\partial N}{\partial x^2}
+\frac{\partial\partial N}{\partial y^2}
=\frac1{nn}\frac{f'''(\xi)f'(\xi)-2(f''(\xi))^2}{(f'(\xi))^2}
=-\frac1{nn}\frac{\theta''(\xi)}{\theta(\xi)},
\tag*{7)}
\end{equation*}
\noindent wo
\[
\theta(\xi)=\frac1{f'(\xi)}=\frac{a\cos\Phi}{\surd(1-ee\sin\Phi^2)}
\]
\noindent ist.

Da hiernach \(\theta'(\xi)=-\sin\Phi\) und
\(\theta''(\xi)=-\cos\Phi\dfrac{(1-ee\sin\Phi^2)^{\frac32}}{a(1-ee)}\) ist, da ferner
\[
r=\frac{a(1-ee)}{\surd(1-ee\sin\Phi^2)^3},\qquad
r'=\frac{a}{\surd(1-ee\sin\Phi^2)},
\]
\noindent so wird
\begin{equation*}
\frac{\partial\partial N}{\partial x^2}
+\frac{\partial\partial N}{\partial y^2}=\frac1{rr'nn}.
\tag*{8)}
\end{equation*}
\noindent Vgl.\ Band VIII, S.\ 385.
\begin{center}\rule{.24\linewidth}{.4pt}\end{center}

Zu der in den Art.\ [15] und [16] gegebenen Azimutreduktion ist noch Folgendes zu
bemerken (vgl.\ den Brief an \textsc{Schumacher} vom 25.\ Juni 1831).
\textsc{Gauss} versteht unter dem Azimut in plano des Punktes \(x_2,y_2\) im Punkt
\(x_1,y_1\) den Winkel, den die beide Punkte verbindende gerade Linie mit der durch
\(x_1,y_1\) gezogenen Parallelen zur \(x\)-Achse bildet. Das Azimut auf dem Sphäroid ist
dagegen der Winkel, den die geodätische Linie in dem Punkt, der \(x_1,y_1\) entspricht,
mit derjenigen Kurve bildet, deren Darstellung in der Ebene eine Parallele zur
Abszissenachse ist. Ist \(x\) nach Süden positiv und bezeichnet im Punkt \((x_1y_1)\)
\(T\) das astronomische Azimut, \(t\) das Azimut auf dem Sphäroid, \(\theta\) das
Azimut in plano, \(c\) die Meridiankonvergenz und \(\psi\) den Winkel, den das Bild
der geodätischen Linie mit der Geraden durch \((x_1y_1)\) und \((x_2y_2)\) bildet, so ist
\[
T=t-c,\qquad \theta=t-\psi.
\]
Die unter [15] aufgeführten Formeln zur Reduktion des Azimuts auf dem Sphäroid auf
das Azimut in plano ergeben sich aus der in Art.\ [16], S.\ 166, abgeleiteten
Gleichung für \(\psi\), wenn darin die Glieder zweiter Ordnung vernachlässigt werden.
Dann ist
\[
\psi=\tfrac12 B_{\frac13}\,.\,(x_2-x_1)-\tfrac12 A_{\frac13}\,.\,(y_2-y_1).
\]
Die Koordinaten \(x_1,y_1\) entsprechen dem Anfangspunkt und die Koordinaten \(x_2,y_2\)
dem Endpunkt der geodätischen Linie. Um die Koeffizienten \(A_{\frac13}\) und
\(B_{\frac13}\) zu erhalten, ist \(\log n\) als Funktion der Koordinaten darzustellen.
Es sei \(x,y\) ein beliebiger Punkt, dem auf dem Ellipsoid ein Punkt mit der Breite
\(\varphi\) entspreche. Setzt man dann
\begin{equation*}
\tfrac12\frac{(1-ee\sin\varphi^2)^2}{aa(1-ee)}=\alpha
\tag*{1)}
\end{equation*}
\par}
\noindent\textsc{ix.}\hfill 26\hspace*{.12\linewidth}\par
\endgroup
